\chapter{Binary search trees}\label{chap:binary:search:trees}

Binary Search Trees (BSTs) are a fundamental data structure for storing sorted data. They allow for efficient insertion, deletion, and lookup operations.

\section{Design rationale}

Our BST implementation uses a standard node-based structure where each node contains a key, a value (for maps), and pointers to left and right children.

To ensure performance, we aim to keep the tree balanced. While a basic BST does not guarantee balance (leading to $O(N)$ worst-case performance), variants like Splay trees or Treaps provide better guarantees.

\section{Binary search tree interface}

The interface provides standard operations, as demonstrated by the Treap interface below:

\begin{code}
  \lstinputlisting[style=C,linerange={3-12}]{libellul/include/libellul/type/treap/interface.h}
  \caption{The generic treap interface.}
  \label{code:treap:interface}
\end{code}

\section{Example variant: Treaps}

Treaps (Tree Heaps) are a randomized binary search tree. Each node is assigned a random priority. The tree maintains the BST property on keys and the Heap property on priorities. This structure tends to be balanced with high probability, offering $O(\log N)$ expected time complexity for all operations.

\section{Example variant: Splay trees}

Splay trees are a self-adjusting binary search tree. When a node is accessed, it is "splayed" (moved) to the root using a series of tree rotations. This property ensures that recently accessed elements are near the root, making them quick to access again. This is particularly useful for applications with temporal locality.
